\documentclass[11pt]{article}
\usepackage[margin=0.85in]{geometry}
\usepackage{amsmath,hyperref}
\title{Attributing methane growth with the 2019 budget}
\author{Reading note on He et al. (2026)}
\date{}
\begin{document}\maketitle
\noindent\textit{Attributing 2019--2024 methane growth using TROPOMI satellite observations.}
Science Advances 12, eadz9007.
\href{https://doi.org/10.1126/sciadv.adz9007}{doi:10.1126/sciadv.adz9007}.
The published article and supplementary Figs.~S1--S9 were obtained.

\paragraph{Method.}
The authors assimilate 610 million bias-corrected TROPOMI retrievals
over 2019--2024 with GEOS-Chem 14.4.1 at $2^\circ\times2.5^\circ$
resolution, driven by GEOS-FP meteorology. Each annual analytical
inversion optimizes 4111 land surface fluxes and two hemispheric OH
scaling factors. The preceding posterior becomes the next prior, but
prior error variances are reset; 2019 is inverted twice as a burn-in.
The base prior errors are 50\% for fluxes and 10\% for OH, with
regularization 0.02. A 27-member ensemble varies these assumptions.
The prior combines EDGARv7 and regional anthropogenic inventories,
LPJ-MERRA2 wetlands, and MeMo soil uptake. Sector attribution uses
the prior sector fractions within each grid cell, rather than independent
sector-specific observation constraints.

\paragraph{Budget and findings.}
Their equation (1) approximates annual methane mass growth as
\begin{equation}
G_j=(G_0-k_0\Delta m_j)+\Delta E_j-m_0\Delta k_{\mathrm{OH},j},
\end{equation}
in which 0 denotes 2019, $m$ is mass, $E$ is gross emissions, and
$k_0$ includes OH and other chemical loss. Soil uptake is a separate
baseline loss contained in $G_0$. Changes in other chemical loss,
soil uptake, and the product of burden and loss-frequency anomalies
are neglected. Emissions rise from 571 Tg yr$^{-1}$ in 2019 to
601 in 2021, then return to 571--575 in 2022--2024. For 2020--2021,
81\% of the forcing beyond the baseline relaxation is assigned to
emissions and 19\% to OH. The 59/25/16\% relaxation/emissions/OH
partition applies to 2019--2024 and must not be relabeled as a
2019--2022 result. Mean OH lifetime is 11.1 years; OH decreases
about 2\% from 2019 to 2022. The largest 2019--2024 regional trends
are East Africa, South America, and Europe.

\paragraph{Interpretation and reproducibility.}
The supplement documents fit errors, averaging kernels, sector overlap,
error correlations, and sensitivity to regularization, but contains no
annual numerical posterior table. The availability statement links model
code and satellite inputs, not the numerical posterior products needed
to redraw maps or diagnostic matrices. Those inputs do not remove the
need for an inversion, which is excluded here. Figure digitization, if
used, must be identified as approximate and kept distinct from original
numerical data. The livestock/wetland attribution is particularly sensitive
to prior spatial overlap and missing tropical observations. The authors
note that precipitation-driven wetland maps can miss river inundation;
their African ammonia comparison does not support rising livestock.
Our decision is to implement their annual decomposition exactly, expose
its neglected terms, and distinguish it from a fully integrated
constant-2019 counterfactual, whose burden differs from the observed one.
\end{document}
